\subsubsection*{Mixed-integer pool expansion}
\label{sec:mip_expansion_formulation}

\emph{Donor attributes, not donor rows.} For a zone $z$ let
$\mathcal{Z}_z^0$ be the set of control cells with a positive census target and
zero attainable capacity in the zone's candidate pool. For each such cell the
module draws donor records from the census microdata that carry the cell. It does
not append the donor person. It takes the donor's person-level attribute vector
as a target and pairs it with a real shell household from the zone's
\texttt{REGUCP}. The donor's \texttt{RLHP} is a preference: if the shell has a
member in that role the solver uses that slot, otherwise the shell slot supplies
the role. The shell supplies the dwelling and family fields and the slot, while
the donor supplies the values to be approximated. A small completion program
constructs the new person as the closest rule-feasible profile to the donor
vector,
\begin{equation}
\min_{y}\;\sum_{a}\bigl(1-y_{a,\tau(a)}\bigr)
\quad\text{s.t.}\quad
\sum_{c} y_{a,c}=1,\;\;
y_{a_1,c_1}+y_{a_2,c_2}\le 1 \;\;\forall\,(a_1,c_1,a_2,c_2)\in\mathcal{C},
\label{eq:mip_repair}
\end{equation}
where $\tau(a)$ is the donor's code for attribute $a$ and $\mathcal{C}$ is the
set of forbidden category pairs obtained by evaluating the shared rule set over
every attribute domain mesh (2{,}027 cuts in the ten-zone development audit).
When the zero cell is a person-level attribute, it is imposed as a hard row
$y_{a^*,c^*}=1$; when it is a household-level field, the shell must already
carry it. When several local shells can host the donor, the constructor prefers a
household size close to the zone's \texttt{nprd} target mix; a size that is the
only carrier for a required cell is still admitted. The solved profile is
written into the shell slot, so the donor's attributes are mixed with the shell's
household context rather than copied as a row. The donor profile may come from the wider population, which is the point of
expansion: it supplies a cell the SA4 sample lacks while the shell stays local.
Every constructed record is tagged to the target SA2 with a fresh identifier, so
sibling SA2s do not share it.

\emph{Donor and shell menus.} The donor source is the entire census, not a
pre-filtered subset. For a zero cell $c$ let $D_c$ be the set of census records
carrying $c$, and let $S_z$ be the local shell households in zone $z$ that can
host the required slot. A menu cap $m_c$ bounds how many donors are sampled per
zero cell, and a shell cap bounds the shells tried per donor; the released
builds use $(m_c,s_c)=(300,32)$ and $(2000,32)$. The larger menu is therefore a
denser sample of the same population, not a different donor source. Every
donor--shell pair still passes the repair and whole-household validation before
it enters $\mathcal{J}_z$, so increasing the cap raises the number of
rule-feasible candidates without relaxing the constraints.

\emph{Family reconciliation.} The household-level fields are fixed to the
shell. After the person-level profile is written, the family and household
counters are recomputed from the solved member vectors: \texttt{CPRF} (persons
in family), \texttt{CACF} (children under 15 in family),
\texttt{persons\_in\_family}, \texttt{persons\_in\_household} and
\texttt{families\_in\_household}. The new person's \texttt{RLHP} role is the
shell slot, and its family and dwelling attributes are those of the block. The
whole household is then validated against the eleven shared rules of
Appendix~\ref{sec:ref}; a candidate that fails is discarded, not patched. This
is what keeps the individual attributes aligned with the household variables: a
donation cannot leave a person with a family role, a child count or a dwelling
term that the household block does not support. Figure~\ref{fig:donor_chromosome}
traces the two steps.

% Figure: donor attributes enter a shell slot, then the family/household
% counters are recomputed and the whole block is validated.
\begin{figure}[t]
\centering
\resizebox{\linewidth}{!}{%
\begin{tikzpicture}[
  >=stealth,
  attr/.style={draw=black!45, minimum width=0.70cm, minimum height=0.42cm,
               inner sep=0pt, font=\scriptsize},
  donor/.style={attr, fill=orange!22, draw=orange!70!black},
  shell/.style={attr, fill=blue!13, draw=blue!55!black},
  repaired/.style={attr, fill=green!22, draw=green!55!black},
  target/.style={attr, fill=orange!50, draw=orange!85!black, line width=0.9pt},
  lab/.style={font=\tiny, anchor=east, inner sep=1pt},
  top/.style={font=\tiny, rotate=90, anchor=west, inner sep=1pt},
  note/.style={font=\tiny, align=left, inner sep=0pt},
  head/.style={font=\scriptsize\bfseries, anchor=west, inner sep=0pt},
  box/.style={draw=black!45, rounded corners=2pt, inner sep=2pt,
              font=\tiny, align=center},
  flow/.style={->, draw=black!60, line width=0.6pt},
  flowd/.style={->, draw=orange!75!black, line width=0.7pt},
  flows/.style={->, draw=blue!60!black, line width=0.7pt},
]

% ======================================================== (a) attribute vector
\node[head] at (0.00,4.70) {(a) donor attributes into a shell slot};

% donor strip
\node[lab] at (0.10,3.40) {donor person};
\foreach \name/\x/\style in {%
  AGEP/1.90/donor, SEXP/2.65/donor, MSTP/3.40/donor,
  HEAP/4.15/donor, QALFP/4.90/target, RLHP/5.65/donor} {
  \node[\style] at (\x,3.40) {};
  \node[top] at (\x-0.06,3.64) {\name};
}
\node[note, anchor=west] at (6.20,3.40) {target cell:\\ \texttt{qalfp:12}};

% shell strip
\node[lab] at (0.10,2.20) {shell household};
\foreach \name/\x in {NPRD/1.90, HIND/2.85, VEHRD/3.80,
                      STRD/4.75, RLHP/5.70} {
  \node[shell, minimum width=0.95cm, font=\tiny] at (\x,2.20) {\name};
}


% new person strip
\node[lab] at (0.10,1.00) {new person};
\foreach \name/\x/\style in {%
  AGEP/1.90/repaired, SEXP/2.65/donor, MSTP/3.40/repaired,
  HEAP/4.15/donor, QALFP/4.90/target, RLHP/5.65/shell} {
  \node[\style] at (\x,1.00) {};
  \node[top] at (\x-0.06,1.24) {\name};
}
% Below the slot it describes, clear of the repaired-profile arrow.
\node[note, anchor=west] at (6.15,0.57)
  {\texttt{qalfp=12, rlhp=3}};

% flows
\draw[flowd] (1.25,3.18) -- (1.25,1.24)
  node[pos=0.25, left, font=\tiny, align=right] {donor\\values};
\draw[flows] (6.35,1.98) -- (6.35,1.24)
  node[pos=0.25, right, font=\tiny] {fixed context};

% repair box
\node[box, text width=2.35cm] (repair) at (7.90,2.60)
  {repair: nearest\\ feasible profile;\\ cuts $\mathcal{C}$};
\draw[flow] (6.40,2.60) -- (repair.west);
% Enter the slot from its right-hand edge, below the "fixed context" arrow,
% so the elbow does not cut through that label.
\draw[flow] (repair.south) |- (6.00,1.10);

% legend
\node[donor, minimum width=0.18cm, minimum height=0.18cm] at (0.20,0.12) {};
\node[note, anchor=west] at (0.32,0.12) {donor};
\node[shell, minimum width=0.18cm, minimum height=0.18cm] at (1.35,0.12) {};
\node[note, anchor=west] at (1.47,0.12) {shell/fixed};
\node[repaired, minimum width=0.18cm, minimum height=0.18cm] at (2.90,0.12) {};
\node[note, anchor=west] at (3.02,0.12) {repaired by the rules};

% ======================================================== (b) family
\node[head] at (9.95,4.70) {(b) family and household reconciliation};

\node[draw=black!45, rounded corners=3pt, fill=blue!4,
      minimum width=5.25cm, minimum height=2.55cm] (hh) at (12.50,2.65) {};
\node[note, anchor=north] at (12.50,3.78)
  {household fields fixed; counters recomputed};

\node[circle, draw=black!55, fill=blue!12, minimum size=0.62cm,
      font=\tiny] (a1) at (10.70,3.15) {1};
\node[circle, draw=black!55, fill=blue!12, minimum size=0.62cm,
      font=\tiny] (a2) at (11.65,3.15) {2};
\node[circle, draw=orange!80!black, fill=orange!45, minimum size=0.62cm,
      font=\tiny, line width=0.9pt] (ch) at (12.60,3.15) {3};
\node[circle, draw=black!55, fill=blue!12, minimum size=0.62cm,
      font=\tiny] (o) at (13.55,3.15) {...};

\node[note, anchor=north] at (10.70,2.80) {\texttt{RLHP=1}};
\node[note, anchor=north] at (11.65,2.80) {\texttt{RLHP=2}};
\node[note, anchor=north] at (12.60,2.80) {\texttt{RLHP=3}\\new};
\node[note, anchor=north] at (13.55,2.80) {other};

\node[box] (cprf) at (10.70,1.95) {\texttt{CPRF}};
\node[box] (cacf) at (11.55,1.95) {\texttt{CACF}};
\node[box] (nprd) at (12.40,1.95) {\texttt{NPRD}};
\node[box] (fam)  at (13.35,1.95) {\texttt{families}};
\draw[flow] (cprf.north) -- ++(0,0.22);
\draw[flow] (cacf.north) -- ++(0,0.22);
\draw[flow] (nprd.north) -- ++(0,0.22);
\draw[flow] (fam.north) -- ++(0,0.22);

% validation
\node[diamond, draw=black!55, fill=blue!5, aspect=2.0, inner sep=1pt,
      font=\tiny, align=center] (val) at (12.50,0.72)
  {11 shared\\rules};
\draw[flow] (12.50,1.55) -- (val.north);
% The outcomes are boxes, so both arrows terminate at a node.
\node[box, draw=red!65!black, fill=red!8, text=red!65!black,
      minimum width=0.80cm] (discard) at (10.60,0.72) {discard};
\node[box, draw=green!55!black, fill=green!9, text=green!45!black,
      minimum width=0.70cm] (emit) at (14.30,0.72) {emit};
\draw[flow, draw=green!55!black] (val.east) -- (emit.west)
  node[midway, above=1.5pt, font=\tiny, text=green!45!black] {pass};
\draw[flow, draw=red!65!black] (val.west) -- (discard.east)
  node[midway, above=1.5pt, font=\tiny, text=red!65!black] {fail};

\node[note, anchor=west] at (9.95,0.05)
  {the rule check is on the whole household, after the counters are recomputed};

\end{tikzpicture}%
}
\caption{Donor attributes are mixed into a shell slot, not copied as a person
row. \textbf{(a)} The donor's person-level vector (orange) is the target; the
shell household (blue) fixes the dwelling fields and supplies the
\texttt{RLHP} slot. The donor's \texttt{RLHP} is a preference: the solver uses a
shell slot with that role when one exists, otherwise the shell slot defines the
new person's role. The completion program returns the closest rule-feasible
profile to the donor vector, with the target cell (\texttt{qalfp:12}) as a hard
row; attributes that the rules force away from the donor value are shown in
green. The solved profile is written into the shell's slot, so the new person
carries the donor's attributes where the rules allow and the shell's household
context everywhere else. \textbf{(b)} The family and household counters
(\texttt{CPRF}, \texttt{CACF}, \texttt{NPRD}, \texttt{families}) are recomputed
from the solved member vectors, and the whole household is checked against the
eleven shared rules of Appendix~\ref{sec:ref}. A candidate that fails is
discarded, not patched; a candidate that passes is emitted with a fresh
identifier and the target SA2 tag. Profiles outside the person-level model fall
back to a whole-household copy; the figure shows the person-level repair path.}
\label{fig:donor_chromosome}
\end{figure}


\emph{Constraint guarantee.} Every candidate is checked as a whole household,
after the counters have been recomputed, against the same eleven rules used by
the pool builder (Appendix~\ref{sec:ref}). The repair program also carries the
forbidden-pair cuts $\mathcal{C}$, so a donor attribute cannot create a
combination the rules forbid; if no feasible profile exists, the donor--shell
pair is rejected. Equation~\eqref{eq:mip_repair} is therefore a
\emph{candidate repair}, not the pool objective: it preserves a real donor
profile wherever the hard rules allow and is solved separately for each
candidate before any selection. For attributes outside the person-level model
(e.g.\ \texttt{tend}, \texttt{typp}, \texttt{indp}, \texttt{stup},
\texttt{assnp}), a person-level repair is not available, so the module copies a
whole real donor household verbatim, recomputes and validates it, and skips it
if it fails. No emitted record leaves the rule set.

\emph{Pool-support block.} The candidates are considered together, not
cell by cell. Let $\mathcal{J}_z$ index the candidate menu,
$y_j\in\{0,1\}$ indicate admission to the zone's candidate set, and $s_j$ be
the number of persons in candidate $j$. For a control cell $(k,m)$, let $a_{j,k,m}$ be the candidate's
contribution, the number of its members in the cell for a person-level table, or
one if the household sits in the cell for a household-level table, and let
$u_{j,k}=s_j$ for a person-level table and $u_{j,k}=1$ for a household-level
table. With the zone target $E_{z,k,m}$ and target share
\begin{equation}
q_{k,m}=\frac{E_{z,k,m}}{\sum_{m'\in\mathcal{M}_k}E_{z,k,m'}},
\label{eq:mip_share}
\end{equation}
each zero-capacity cell $c=(k,m)\in\mathcal{Z}_z^0$ has a carrier requirement
\begin{equation}
\nu_c \;=\; \left\lceil \frac{E_{z,k,m}}{r_{\max}} \right\rceil ,
\label{eq:mip_carrier_requirement}
\end{equation}
the number of carrier contribution units that must be selected so that
$r_{\max}$ replicas of each selected household can reach its target. The
support block is enabled by default. When it is on, the following constraints
define the admitted set used by the zone program below and the admitted pool is
\(\mathcal{J}_z^{+}=\{j\in\mathcal{J}_z:y_j=1\}\); when it is off, the block is
omitted and \(\mathcal{J}_z^{+}=\mathcal{H}_z\), so the fit has no access to
any constructed record. The block adds no objective term.
\begin{equation}
\begin{aligned}
& \sum_{j\in\mathcal{J}_z} a_{j,c}\,y_j+e_c\ge \nu_c,
&& c\in\mathcal{Z}_z^0,\\
& d_{k,m}\ge \sum_{j\in\mathcal{J}_z}
  \bigl(a_{j,k,m}-q_{k,m}u_{j,k}\bigr)y_j,
&& k\in\mathcal{K},\,m\in\mathcal{M}_k,\\
& d_{k,m}\ge -\sum_{j\in\mathcal{J}_z}
  \bigl(a_{j,k,m}-q_{k,m}u_{j,k}\bigr)y_j,
&& k\in\mathcal{K},\,m\in\mathcal{M}_k,\\
& \sum_{j\in\mathcal{J}_z}y_j\le K_z,\qquad
  \sum_{j\in\mathcal{J}_z}s_jy_j\le P_z^{\mathrm{add}},\\
& \sum_{j\in\mathcal{J}_z(h)}y_j\le 1,
&& h\in\mathcal{H}_z,\\
& y_j\in\{0,1\},\quad e_c\ge 0,\quad d_{k,m}\ge 0.
\end{aligned}
\label{eq:mip_joint}
\end{equation}
The support variable \(y_j\) governs admission to the candidate set, not use in
the fitted population: \(y_j=1\) makes candidate \(j\) available to the zone
fitter, which then selects from that set with \(x_{i,h}\). The slack \(e_c\)
absorbs any shortfall below the carrier requirement, and \(d_{k,m}\) records the
admitted set's deviation from the zone's marginal shares; neither appears in
\eqref{eq:objective}. Among feasible support sets the implementation prefers
lower coverage slack, then lower marginal deviation, then fewer admitted
households, and scalarises that preference internally as
$\lambda_0\sum_c e_c+\sum_{k,m}w_{k,m}d_{k,m}+\lambda_1\sum_j y_j$, with
$\lambda_0$ large enough that coverage dominates. That scalarisation is a
construction tie-break inside the block, not a second minimisation of
the fitted population, and it does not compete with RSSZ. Here
$w_{k,m}=1/(|\mathcal{M}_k|\max(E_{z,k,m},1))$ gives every control table equal
weight; the released defaults are $\lambda_0=10^5$ and $\lambda_1=10^{-2}$.
Coverage remains soft, so an unsourceable cell is reported rather than making
the block infeasible.

With the block on, the zone program below minimises \eqref{eq:objective} over
\(\mathcal{J}_z^{+}\); with it off, the same program minimises \eqref{eq:objective}
over \(\mathcal{H}_z\). The fitter, the objective and constraints C1--C7 are
unchanged in the two runs, so the difference in fitted RSSZ is attributable to
the extra pool support, not to a different fitting criterion. This is the
counterfactual the expanded-pool arms are designed to measure.

$K_z$ and $P_z^{\mathrm{add}}$ are not fixed budgets but are derived from the
zero cells and their targets:
\begin{equation}
K_z=\sum_{c\in\mathcal{Z}_z^0}\nu_c,
\qquad
P_z^{\mathrm{add}}=M K_z,
\label{eq:mip_budget}
\end{equation}
where $M$ is the household-size cap. The budget therefore scales with the
marginal gap and cannot bind below what coverage requires. The constraint over
$\mathcal{J}_z(h)\subseteq\mathcal{J}_z$ permits at most one transplant per
shell household.

\emph{Pool growth.} The program above determines the coverage-minimal set. The
constructed menu contains at least $\nu_c$ donor variants for each zero cell, so
the selection is not limited by a fixed menu breadth. The released
implementation can also append the rest of the constructed candidate
menu for the zone, not only the selected households. Those extra records are
tagged to the same SA2 and receive fresh identifiers, so the SA2 fitter has more
candidates to choose from while no constructed record is shared with a sibling
SA2. The coverage inequality is therefore a lower bound on the addition, not a
ceiling on the pool.

The important difference from repairing cells one at a time is that
\eqref{eq:mip_joint} uses \emph{one} admission vector $y$ for \emph{all} control
tables. The deviations $d_{k,m}$ are therefore not independent cell repairs:
they measure how far the selected agents' aggregate profile is from the zone's
total marginal composition, table by table, and the solver must trade those
deviations against one another and against coverage. A set of records that fills
every zero cell but distorts the income or tenure shares is penalised; a set that
matches the shares but leaves a zero cell uncovered is penalised more heavily
through $\lambda_0$. This is a joint fit over the marginal tables, not a fit to
the full joint distribution: one selection vector must satisfy every table
simultaneously, while the joint distribution over attribute combinations is the
object of the zone fitter, not of the small pool-expansion program. The
coverage constraints are soft ($e_c$) precisely so that an unsourceable
cell, a structural zero, a padding code, or a cell with no donor anywhere, does
not make the program infeasible; it is reported instead.

The model is deliberately \emph{not} required to reproduce the census at pool
scale. The deficit in a zone is often tens of thousands of persons, while the
candidate pool is a 5\% microsample. Requiring
$\sum_j a_{j,k,m}y_j\ge E_{z,k,m}$ at pool construction would force the pool to
carry the census many times over, which is the fitter's job, not the pool's.
The construction requirement is instead on \emph{carrier contributions}:
$\sum_j a_{j,c}y_j\ge \nu_c=\lceil E_{z,k,m}/r_{\max}\rceil$ for every
zero-capacity cell, so that $r_{\max}$ replicas of each selected household can
reach the cell target. This is weaker than $E_{z,k,m}$ by the replication
factor $r_{\max}$ and sufficient for coverage; partial deficits outside the
zero-cell set remain the fitter's job. Because coverage is soft, a cell can
still be left below $\nu_c$ when the candidate menu makes the requirement
impossible; such cells are reported rather than silently dropped. The zone fit
that follows decides how often each selected candidate is actually used.
